% 4-4-direct-vs-agg.tex
%  Budget: 2pp.  RQ4. HANDOFF 6.2.
%  Source map: no external claims; this is an internal comparison of two arms
%   built on the same feature matrix, origins and tuning budget.
%  Values from reports/tables/results_canonical.csv (MAE):
%   SARIMAX   direct 3.1976  agg 3.2691  -> DIRECT wins (the exception)
%   LEAR      direct 2.8989  agg 2.8389  -> aggregation
%   LightGBM  direct 3.3012  agg 2.9928  -> aggregation
%   LSTM      direct 3.1753  agg 2.7804  -> aggregation
%   naive     identical both ways
%  MANDATE: explain the SARIMAX exception, do not average it away.

\section{مقایسه روش مستقیم و تجمیعی}
\label{sec:4-4}

پرسش پژوهشی چهارم مستقیم است: برای پیش‌بینی میانگین بار پایه‌ی روزانه، کدام مسیر دقت
بالاتری دارد \lr{—} مدل‌سازی مستقیمِ هدف روزانه، یا میانگین‌گیری از ۲۴ پیش‌بینی
ساعتی؟ و آیا پاسخ برای همه‌ی مدل‌ها یکسان است؟

پاسخ به بخش دوم پرسش منفی است، و همین آن را به یافته بدل می‌کند.

\begin{table}[htbp]
\centering
\caption{مقایسه‌ی مستقیم دو مسیرِ رسیدن به هدف روزانه بر معیار \lr{MAE}. ستون
«اختلاف» مقدار مسیر تجمیعی منهای مسیر مستقیم است؛ مقدار منفی یعنی تجمیع بهتر بوده
است. مدل ساده در هر دو مسیر یکسان است.}
\label{tab:direct-vs-agg}
\small
\setlength{\tabcolsep}{6pt}
\begin{tabular}{lrrrl}
\toprule
مدل & مستقیم & تجمیعی & اختلاف & برنده \\
\midrule
\lr{SARIMAX} & \lr{\textbf{3.1976}} & \lr{3.2691} & \lr{+0.0715} & مستقیم \\
\lr{LEAR-LASSO} & \lr{2.8989} & \lr{\textbf{2.8389}} & \lr{-0.0600} & تجمیعی \\
\lr{LightGBM} & \lr{3.3012} & \lr{\textbf{2.9928}} & \lr{-0.3084} & تجمیعی \\
\lr{LSTM} & \lr{3.1753} & \lr{\textbf{2.7804}} & \lr{-0.3949} & تجمیعی \\
\midrule
ساده & \lr{6.3575} & \lr{6.3575} & \lr{0.0000} & یکسان \\
\bottomrule
\end{tabular}
\end{table}

\subsection*{تجمیع در سه مورد از چهار مورد برنده است}

برای \lr{LEAR-LASSO}، \lr{LightGBM} و \lr{LSTM}، میانگین‌گیری از ۲۴ پیش‌بینی ساعتی
دقیق‌تر از مدل‌سازی مستقیم است. اندازه‌ی این برتری یکسان نیست: برای \lr{LEAR-LASSO}
تنها \lr{0.0600} است، اما برای \lr{LightGBM} \lr{0.3084} و برای \lr{LSTM}
\lr{0.3949}.

توضیح این الگو در مقدار اطلاعاتی است که هر مسیر در اختیار مدل می‌گذارد. مسیر تجمیعی،
مدل را وادار می‌کند ساختار درون‌روزی را یاد بگیرد \lr{—} یعنی ۲۴ هدف مجزا با ۲۴
الگوی متفاوت \lr{—} و سپس آن ۲۴ پیش‌بینی را میانگین می‌گیرد. خطاهای ساعت‌های مختلف
کاملاً هم‌بسته نیستند، پس بخشی از آن‌ها در میانگین‌گیری یکدیگر را خنثی می‌کنند. مسیر
مستقیم، در مقابل، این ساختار را از همان ابتدا دور می‌ریزد و تنها یک عدد در روز را
هدف می‌گیرد.

اینکه بیشترین سود از آنِ دو مدل غیرخطی است نیز با همین توضیح سازگار است: مدل‌های
انعطاف‌پذیرتر توانِ بیشتری برای بهره‌بردن از ۲۴ الگوی مجزا دارند، و در مقابل، وقتی به
یک هدفِ تک‌بعدی محدود شوند بیشترین سهم از ظرفیتشان بلااستفاده می‌ماند.

\subsection*{\lr{SARIMAX} استثناست، و این استثنا واقعی است}

\lr{SARIMAX} تنها مدلی است که مسیر مستقیم برایش بهتر عمل می‌کند: \lr{3.1976} در
برابر \lr{3.2691}. این نتیجه را نباید به‌عنوان نوفه کنار گذاشت یا در میانگینی پنهان
کرد؛ توضیحی ساختاری دارد.

\lr{SARIMAX} یک مدل کلاسیک سری زمانی است که برای سری‌های تک‌بعدی با ساختار
خودهمبستگی و مؤلفه‌ی فصلی طراحی شده است. سری بار پایه‌ی روزانه دقیقاً چنین سری‌ای است:
هموار، تک‌بعدی، با تناوب هفتگی روشن. وقتی این مدل مستقیماً بر آن سری برازش می‌شود، در
زیستگاه طبیعی خودش قرار دارد.

در مسیر تجمیعی، همین مدل باید ۲۴ سریِ ساعتیِ به‌مراتب پرنوفه‌تر را مدل کند و سپس
نتایج را میانگین بگیرد. هر یک از آن ۲۴ برازش، خطای خود را دارد، و برخلاف مدل‌های
انعطاف‌پذیرتر، \lr{SARIMAX} نمی‌تواند از ساختار مشترک میان ساعت‌ها بهره ببرد، چون
۲۴ مدل کاملاً مستقل‌اند. افزون بر آن، همان‌گونه که \cref{sec:3-6} ثبت کرد، این مدل تنها هر
هفت روز به‌طور کامل بازبرازش می‌شود، و این محدودیت در ۲۴ سریِ پرنوفه هزینه‌ی بیشتری
دارد تا در یک سریِ هموار.

\subsection*{چرا این مقایسه منصفانه است}

اعتبار این بخش تماماً به یک نکته بستگی دارد: هر دو مسیر باید تلاش برابری دریافت کرده
باشند. اگر مسیر مستقیم بودجه‌ی تنظیم کمتری می‌گرفت، هر برتریِ مسیر تجمیعی می‌توانست
صرفاً بازتاب همان نابرابری باشد و نه خاصیتِ خودِ مسیر.

سه چیز میان دو مسیر یکسان نگه داشته شده است: همان ماتریس ویژگی، همان مبادی پیش‌بینی،
و همان بودجه‌ی تنظیمِ پنجاه‌آزمایشی به‌ازای هر مدل. مسیر مستقیم پیکربندی تنظیم‌شده‌ی
مستقل خود را دارد، که با همان شمار آزمایش و بر همان پنجره‌ی اعتبارسنجی به دست آمده
است. این برابری، طراحیِ آگاهانه بوده و همان چیزی است که نتیجه‌ی این بخش را از یک مصنوعِ
طراحی جدا می‌کند.

\subsection*{آنچه این بخش ادعا نمی‌کند}

هیچ قاعده‌ی کلی درباره‌ی تجمیع زمانی در پیش‌بینی از اینجا به دست نمی‌آید. چهار مدل، یک
بازار، یک پنجره‌ی دوساله. آنچه می‌گوید این است که در این شرایط، تجمیع برای سه مدل از
چهار مدل بهتر بود و استثنای چهارم توضیح ساختاری دارد \lr{—} و اینکه پرسش پژوهشی
عامدانه چنان صورت‌بندی شده بود که این استثنا بخشی از پاسخ باشد و نه دردسری برای آن.
